% ============================================================
% Plantilla APA 7.ª edición (modo estudiante) con LaTeX
% Compilar:  pdflatex -> biber -> pdflatex -> pdflatex
%        o:  latexmk -pdf articulo_apa7.tex
% En Overleaf: Menu > Compiler = pdfLaTeX; biber se ejecuta solo.
% ============================================================
\documentclass[stu,12pt,a4paper,floatsintext]{apa7}
% Modos: stu = estudiante | man = manuscrito | jou = revista | doc = documento
% Para Letter (EE. UU.) use letterpaper en lugar de a4paper.

\usepackage[utf8]{inputenc}
\usepackage[T1]{fontenc}
\usepackage[spanish,es-tabla]{babel}  % es-tabla: «Tabla» en vez de «Cuadro»
\usepackage{csquotes}                        % obligatorio con biblatex + babel
\usepackage{amsmath}
\usepackage{booktabs}
\usepackage{graphicx}
\usepackage{caption}

% --- Bibliografía: biblatex-apa + biber ---------------------
\usepackage[style=apa,backend=biber,sortcites=true,sorting=nyt]{biblatex}
\DeclareLanguageMapping{spanish}{spanish-apa}
\addbibresource{referencias_apa.bib}

% --- Portada (modo estudiante) -------------------------------
\title{Respuesta transitoria de un circuito RC: verificación experimental}
\shorttitle{Circuito RC}   % encabezado de página (máx. 50 caracteres)
\authorsnames{Nombre Apellido}
\authorsaffiliations{Facultad de Ingeniería Electrónica, Universidad Nacional Tecnológica de Lima Sur}
\course{IE-101: Circuitos Eléctricos}
\professor{Dr. Nombre Apellido}
\duedate{27 de octubre de 2026}

\abstract{%
  El resumen APA tiene entre 150 y 250 palabras, en un solo párrafo sin sangría.
  Describa el problema, el método, los resultados principales y la conclusión.
  Este ejemplo reproduce la estructura de un informe breve sobre la carga de un
  capacitor y muestra cómo citar, numerar tablas y figuras, y generar la lista
  de referencias automáticamente.}
\keywords{circuito RC, respuesta transitoria, constante de tiempo, APA 7, LaTeX}

\AtBeginDocument{\renewcommand{\notelabel}{Nota}}

\begin{document}
\maketitle

% apa7 repite el título (negrita, centrado) al inicio del cuerpo.
% Nivel 1 = \section (centrado, negrita); nivel 2 = \subsection; nivel 3 = \subsubsection.

\subsection{Introducción}
El comportamiento transitorio de los circuitos de primer orden es un tema
central en el análisis de circuitos \parencite{boylestad2023}. Según
\textcite{sedra2020}, la constante de tiempo determina la rapidez de respuesta
del sistema. Las normas de formato siguen el manual de la
\textcite{apa2020}, y el trabajo se compuso con el paquete
\textcite{weiss_apa7}.

\subsubsection{Citas de ejemplo}
Una cita \emph{narrativa} integra al autor en la oración:
\textcite{shannon1948} formalizó la teoría de la información. Una cita
\emph{entre paréntesis} lo coloca al final \parencite{kasneci2023}. Con dos
fuentes: \parencite{vaswani2017,shannon1948}. Con página:
\parencite[p.~381]{shannon1948}. Con más de tres autores, biblatex-apa abrevia
automáticamente a «et al.» (\cite{kasneci2023}).

\subsection{Método}
La carga de un capacitor sigue la ecuación~\eqref{eq:carga}:
\begin{equation}
  v_C(t) = V_{in}\left(1 - e^{-t/\tau}\right), \qquad \tau = RC
  \label{eq:carga}
\end{equation}

\subsection{Resultados}
La Tabla~\ref{tab:carga} resume las mediciones y la Figura~\ref{fig:carga}
muestra la curva obtenida.

% APA 7: número en negrita, título en cursiva, sin líneas verticales.
\begin{table}[htbp]
  \caption{Voltaje en el capacitor durante la carga}
  \label{tab:carga}
  \centering
  \begin{threeparttable}
  \begin{tabular}{cc}
    \toprule
    Tiempo (ms) & Voltaje (V) \\
    \midrule
    0  & 0{,}0 \\
    10 & 2{,}1 \\
    20 & 3{,}5 \\
    30 & 4{,}3 \\
    40 & 4{,}7 \\
    50 & 4{,}9 \\
    \bottomrule
  \end{tabular}
  \begin{tablenotes}[para]
    \item \textit{Nota.} $V_{in} = 5$~V, $R = 10$~k$\Omega$, $C = 3{,}3$~\textmu F.
    Datos ilustrativos.
  \end{tablenotes}
  \end{threeparttable}
\end{table}

% Figura sin archivo externo: usa un recuadro para que compile siempre.
\begin{figure}[htbp]
  \caption{Curva de carga del capacitor}
  \label{fig:carga}
  \centering
  \fbox{\parbox{0.7\textwidth}{\centering\vspace{2cm}%
    Reemplace por \texttt{\textbackslash includegraphics\{curva.pdf\}}%
    \vspace{2cm}}}
  \figurenote{Elaboración propia. Si la figura proviene de otra
  fuente, cite: Adaptado de \textit{título}, por Autor, año, editorial.}
\end{figure}

\subsection{Discusión}
Los valores medidos concuerdan con el modelo teórico, con diferencias
atribuibles a la tolerancia de los componentes.

\subsection{Uso de herramientas de IA}
En APA 7, el texto generado por IA se reconoce explícitamente. Ejemplo:
\emph{Para depurar el código \LaTeX\ se consultó ChatGPT}
\parencite{openai_chatgpt}; la verificación y la redacción final son del autor.

% Referencias: se generan automáticamente.
\printbibliography[title={Referencias}]

\end{document}
